%%%%%%%%%%%%%%%%%%%%%%%%%%%%%%%%%%%%%%%%%%%%%%%%%%%%%%%%%%%%%%
%%%% MA-0714 Optimización %%%%%%%%%%%%%%%%%%%%%%%%%%%%%%%%%%%%
%%%%%%%%%%%%%%%%%%%%%%%%%%%%%%%%%%%%%%%%%%%%%%%%%%%%%%%%%%%%%%
\newpage
\subsection*{MA-0714 Optimización}
\addcontentsline{toc}{subsection}{MA-0714 Optimización}
\hypertarget{MA0714}{}
\begin{refsection}
\begin{table}[H]
\centering{
\begin{tabular}{|ll|}
\hline
\textbf{Año:} IV  & \textbf{Requisitos:} Por definir \\
\textbf{Ciclo:} Optativa  & \textbf{Correquisitos:} Ninguno \\
\textbf{Tipo de curso:} Teórico & \textbf{Horas presenciales por semana:}   \\
\textbf{Créditos:} 4 & \textbf{Horas de trabajo independiente por semana:}  \\ \hline
\end{tabular}
}
\end{table}

\subsubsection*{Descripción del curso}

La optimización estudia la existencia y la caracterización de los valores extremos de una función sobre un conjunto factible. Este primer curso presenta sus fundamentos desde una perspectiva propia de la matemática pura: se da prioridad a las definiciones precisas, los enunciados completos, las demostraciones y la relación entre las hipótesis de cada teorema y sus conclusiones.

A partir de herramientas del análisis real en varias variables, el álgebra lineal y la geometría convexa, se estudian de manera articulada los problemas sin restricciones, la programación lineal y los problemas con restricciones de igualdad y desigualdad. Los resultados de existencia, las condiciones necesarias y suficientes de optimalidad, los teoremas de separación y la dualidad constituyen el eje del curso. Los algoritmos clásicos se utilizan únicamente para ilustrar las consecuencias de la teoría y no como sustituto de su desarrollo matemático.

El curso busca que la persona estudiante formule problemas de optimización con precisión, identifique las hipótesis que permiten garantizar o caracterizar sus soluciones y construya demostraciones rigurosas. De esta manera, sirve como base para estudios posteriores en análisis convexo, optimización no lineal, investigación de operaciones y áreas afines.

\subsubsection*{Objetivos}

\textbf{General:} Desarrollar los fundamentos matemáticos de la optimización para establecer la existencia y caracterizar las soluciones de problemas con y sin restricciones.

\textbf{Específicos:} Al finalizar el curso se espera que la persona estudiante sea capaz de:

\begin{enumerate}
\item Formular con precisión problemas de optimización e identificar sus conjuntos factibles y sus soluciones locales y globales.
\item Establecer la existencia de soluciones mediante argumentos topológicos, en particular a partir de la compacidad y la semicontinuidad o continuidad de la función objetivo.
\item Deducir y aplicar condiciones necesarias y suficientes de primer y segundo orden para problemas sin restricciones.
\item Emplear la geometría de los conjuntos convexos, los teoremas de separación y las propiedades de las funciones convexas para demostrar resultados de optimalidad global.
\item Analizar la estructura geométrica de los problemas de programación lineal y demostrar sus resultados fundamentales de dualidad.
\item Deducir las condiciones de los multiplicadores de Lagrange y de Karush--Kuhn--Tucker, distinguiendo el papel de las cualificaciones de restricciones.
\item Establecer condiciones de suficiencia y resultados básicos de dualidad para problemas convexos.
\item Justificar matemáticamente los principios en que se basan algunos algoritmos clásicos, sin convertir su implementación en el objetivo central del curso.
\item Consultar literatura especializada y comunicar argumentos matemáticos de manera clara, rigurosa y efectiva, tanto en forma oral como escrita.
\end{enumerate}

\subsubsection*{Contenidos}
\begin{enumerate}
\item \textbf{Fundamentos y optimización sin restricciones (5 semanas):}
\begin{enumerate}
    \item Formulación general de un problema de optimización. Óptimos locales y globales. Conceptos topológicos básicos.
    \item Existencia de soluciones: compacidad, continuidad y teorema de valores extremos. Nociones elementales de coercividad.
    \item Condiciones necesarias y suficientes de primer y segundo orden para óptimos locales.
    \item Conjuntos convexos, combinaciones convexas y teoremas de separación.
    \item Funciones convexas y estrictamente convexas. Caracterizaciones de primer y segundo orden y optimalidad global.
    \item Optimización cuadrática sin restricciones como ejemplo integrador.
    \item Principios teóricos de los métodos de descenso como ilustración de la teoría.
\end{enumerate}

\item \textbf{Programación lineal (4 semanas):}
\begin{enumerate}
    \item Formas canónica y estándar. Poliedros, puntos extremos, direcciones extremas y bases.
    \item Existencia de soluciones y caracterización de soluciones óptimas mediante puntos extremos.
    \item Lema de Farkas y alternativas lineales.
    \item Problema dual. Dualidad débil, dualidad fuerte y holgura complementaria.
    \item Fundamento geométrico del método símplex; ejemplos y aplicaciones.
\end{enumerate}

\item \textbf{Optimización con restricciones de igualdad y desigualdad (5 semanas):}
\begin{enumerate}
    \item Restricciones de igualdad: espacio tangente y teorema de los multiplicadores de Lagrange.
    \item Restricciones de desigualdad: direcciones factibles, condiciones de Karush--Kuhn--Tucker y papel de las cualificaciones de restricciones.
    \item Condiciones suficientes de optimalidad en problemas convexos.
    \item Funciones cóncavas y cuasicóncavas como formulaciones equivalentes para problemas de maximización.
    \item Función dual de Lagrange, dualidad débil y condiciones para la dualidad fuerte en programación convexa.
    \item Optimización cuadrática con restricciones de igualdad y sistema KKT como ejemplo integrador.
\end{enumerate}

\item \textbf{Proyectos de descubrimiento (2 semanas):}
Temas relacionados con la materia para ser desarrollados por las personas estudiantes mediante la lectura guiada de literatura matemática. Los temas serán seleccionados de común acuerdo con la persona docente; por ejemplo, métodos de punto interior, subgradientes y optimización no diferenciable, dualidad de Fenchel o aplicaciones de la optimización convexa.

\end{enumerate}

\subsubsection*{Metodología}

\noindent \textbf{Lineamientos metodológicos:} La persona docente a cargo de este curso debe respetar las siguientes pautas:
\begin{enumerate}
    \item La presentación de los contenidos debe priorizar el rigor matemático: definiciones precisas, enunciados completos y demostraciones de los resultados fundamentales (existencia, condiciones de optimalidad, suficiencia y dualidad).
    \item Los algoritmos de optimización se abordan principalmente desde la perspectiva de su justificación teórica y como ilustración de los resultados, sin que la implementación computacional sea el eje del curso.
    \item Cada estudiante debe realizar una revisión de la bibliografía como complemento al trabajo presencial.
\end{enumerate}

\textbf{Sugerencias metodológicas:} Adicionalmente, se tienen las siguientes recomendaciones para la persona docente a cargo de este curso:
\begin{enumerate}
    \item Aprovechar los problemas de optimización cuadrática, sin y con restricciones de igualdad, como ejemplos recurrentes que conectan las condiciones de optimalidad con el álgebra lineal.
    \item Utilizar la programación lineal y su dualidad como caso prototípico, de modo que las condiciones de Karush-Kuhn-Tucker y la dualidad de Lagrange aparezcan posteriormente como su generalización natural.
    \item Promover el trabajo constante a lo largo del ciclo lectivo mediante tareas que combinen demostraciones, construcción de ejemplos y contraejemplos, y resolución analítica de problemas concretos.
    \item Incorporar proyectos de descubrimiento sobre temas que extiendan la materia del curso, definidos de común acuerdo con la persona docente.
    \item Fomentar la comunicación clara y efectiva de ideas matemáticas, tanto de forma oral como escrita, sugiriendo el uso de \LaTeX{} para la presentación de resultados.
\end{enumerate}

\subsubsection*{Orientación bibliográfica}

Como introducciones generales a la teoría de optimización se recomiendan \textcite{Sun96} y \textcite{LY16}. La programación lineal, su interpretación geométrica y su dualidad pueden estudiarse con mayor profundidad en \textcite{BT97}. Para convexidad, programación convexa y dualidad, la referencia principal es \textcite{BV04}. Los problemas no lineales, las condiciones de optimalidad, las cualificaciones de restricciones y los fundamentos de los métodos clásicos se complementan con \textcite{Rus06} y \textcite{Ber16}.

\nocite{Sun96}
\nocite{LY16}
\nocite{BT97}
\nocite{BV04}
\nocite{Rus06}
\nocite{Ber16}

\printbibliography[heading=subbibliography,section=\the\value{refsection}]
\end{refsection}
